%% Hu_Ha
Si besoin est, $T_{s}$ est converti en Kelvin, puis :
\[
e = \frac{H_{a}}{1000} R_v T_{s}
\]
\[
H_{u} = 100 \frac{e}{e_{w}(T_{s})}
\]
Avec :\newline
$R_v = 461.51 J kg^{-1} K^{-1}$ (constante sp�cifique de la vapeur d'eau).\newline
La fonction $e_w(T)$ (tension de vapeur saturante au-dessus de l'eau liquide)
est celle de la biblioth�que thermodynamique d'AIDA, d�crite en annexe 
\ref{annexe-thermo-ew}.
